% % % \documentclass[aps, prl, preprint]{revtex4-2} % \documentclass{report}

% % \usepackage{graphicx} % \usepackage{dcolumn} %
% \usepackage{array,multirow,adjustbox} % \usepackage{bm} %
% \usepackage{amsfonts, amssymb, amsmath} % \usepackage{braket} %
% \usepackage{siunitx}

% % % \newcommand{\ket}[1]{\textstyle \left| #1 \right\rangle \displaystyle} % %
% \newcommand{\bra}[1]{\textstyle \left\langle #1 \right| \displaystyle} % %
% \newcommand{\braket}[2]{\textstyle \left\langle #1 \left|  #2
% \right\rangle\right. \displaystyle}

% % \let\originalleft\left % \let\originalright\right %
% \renewcommand{\left}{\mathopen{}\mathclose\bgroup\originalleft} %
% \renewcommand{\right}{\aftergroup\egroup\originalright}

% % \newcommand{\im}{\mathrm{i}} % \newcommand{\e}{\mathrm{e}} %
% \newcommand{\pwrt}[2]{\frac{\partial #1}{\partial #2}} %
% \newcommand{\diff}{\mathrm{d}}

% % \DeclareMathOperator{\atantwo}{atan2}

% % \title{An open source software package to compare and extend effective mass
% models of the phosphorous donor in silicon} % \author{Luke Pendo} % \date{}

% % \begin{document}

\subsection{Kohn-Luttinger and donor states in silicon}

In a 1955 paper\cite{KohnLuttinger1} W. Kohn and J.M Luttinger derived an
expression for how electrons and holes would move in the presence of a localized
perturbation to the periodic field, in particular the periodic fields of the Si
and Ge semiconductors. These semiconductors do not have a simple conduction band
with a single energy minima -- instead they have 6(Si) or 8(Ge) equivalent
minima lying along different directions in $\mathsfbfit{k}$-space. In a later
paper that year\cite{KohnLuttinger2}, Kohn and Luttinger (hereinafter referred
in this section as KL) presented an application of the effective mass theory to
donor states in silicon. An excess electron, bound by an impurity in silicon
with an excess positive charge, becomes trapped in the crystal. Such a system
can be referred to as an electron donor, conventionally just called a donor. In
silicon, the focus is on the Group V substitutional impurities, such as
phosphorous, which feature an outer electron shell that matches the outer shell
of the silicon atom. This allows the phosphorous to participate in the crystal
structure as though it were silicon. The net effect to the system through the
introduction of this impurity is some minor deformation of the lattice at or
near the donor site, and most importantly that such an atom participating in
silicon bonds will have an excess positive charge that isn't satisfied by the
electrons from its crystalline neighbors. This creates the situation above,
where an excess electron is localized in the region of the donor. While deep
donors exist, we limit our focus here to "shallow" donors. We will save further
discussion of "shallowness" and its impact on the definition of the effective
mass approximation in Section \ref{Shallow}.

KL examined a single-valley model of the donor electron wavefunction that went
as so
\begin{equation}
\Psi_{\rvec}^{}
    =
        F_{\rvec}^{}  \e^{\im \kvec_{\q}\cdot\rvec}u_{0\kvec_{\q};\rvec}
\end{equation}
Above $\kvec_{\q}$ is the point in $k$-space where the energy minima of the
lowest conduction band are located. $F$ represents an envelope function that
introduces spatial localization of the wavefunction. We can see that the
wavefunction takes the form of a conduction band wavefunction attenuated and
made a bound state by a leading envelope function factor.

KL did not extend their model to consider intervalley coupling, and so this
model's ground state manifold is 6-fold degenerate. In actuality, intervalley
couplings lift this 6-fold degeneracy to create a 1-fold $A_{1}$ state,
2-fold $E$ state, and a 3-fold $T_{2}$ state, with $A_{1}$, $E$ and $T_{2}$
referring to the particular symmetries of the overall state. For now, we will
table the discussion of multi-valley states. We will return to this topic in the
next section. KL asserted that in the "shallow" limit when the effective mass
approximation can be applied, $F$ will obey the following equation.
\begin{equation}
\left[
        -
    \frac{\hbar^{2}}{2\real{m}_{\parallel}}\frac{\partial^{2}}{\partial\x^{2}}
        -
    \frac{\hbar^{2}}{2\real{m}_{\parallel}}\frac{\partial^{2}}{\partial\y^{2}}
        -
    \frac{\hbar^{2}}{2\real{m}_{\perp}}\frac{\partial^{2}}{\partial\z^{2}}
        -
    \frac{\e^2}{4\pi\epsilon\sqrt{\x^2+\y^2+\z^2}}
\right]
F_{\rvec}^{}
    =
        \real{E} F_{\rvec}^{}
\label{definition:SV-envelope}
\end{equation}
This is a nonseparable partial differential equation, and no closed form exact
solution is known. However, we can make an educated guess as to the form of the
envelope, and KL predicted that the envelope could be approximated as
\begin{equation}
F_{\rvec}^{}
    \approx
F_{\real{a}\real{b};\rvec}^{\prime}
    \equiv
        \frac{1}{\sqrt{\pi \real{a}^{2}\real{b}}}
        \e^{-\sqrt{\frac{\x^2+\y^2}{\real{a}^2}+\frac{\z^{2}}{\real{b}^2}}}.
\label{definition:KL-wavefunction}
\end{equation}
It is convenient to assert the above function is normalized in the
sense that
\begin{equation}
\int\diff\rvec\,\left|F_{\real{a}\real{b};\rvec}^{\prime}\right|^{2} = 1.
\end{equation}
With a trial function such as this, one which includes arbitrary parameters, one
can compute the expectation energy manifold as a function of the arbitrary
parameters. KL formally combined \eqref{definition:SV-envelope} and
\eqref{definition:KL-wavefunction} to create an effective energy function.
\begin{align}
E_{\real{a}\real{b}}
    & = 
        \frac{\hbar^{2}}{2\real{m}_{\parallel}}
        \int\diff\rvec\,
        \tilde{F}_{\real{a}\real{b};\rvec}^{\prime}
        \left[
                -
            \frac{\partial^{2}}{\partial\x^{2}}
                -
            \frac{\partial^{2}}{\partial\y^{2}}
                -
            \real{\gamma}\frac{\partial^{2}}{\partial\z^{2}}
                -
            \frac{\e^2}{4\pi\epsilon\real{r}}
        \right]
        F_{\real{a}\real{b};\rvec}^{\prime}
,\\
    & =
        \frac{\mathfrak{E}\mathfrak{r}^{2}}{3\real{a}^{2}}
        \left[
            2
                +
            \real{\gamma}\left(\frac{\real{a}}{\real{b}}\right)^2
                -
            6
            \left(\frac{\real{a}}{\real{b}}\right)
            \left(\frac{\real{a}}{\mathfrak{r}}\right)
            \frac
            {
                \text{ArcTan}\left(
                    \sqrt{\left(\frac{\real{a}}{\real{b}}\right)^{2}-1}
                \right)
            }
            { \sqrt{\left(\frac{\real{a}}{\real{b}}\right)^{2}-1} }
        \right]
\label{equation:KL-energy}
.
\end{align}

The quantities $\mathfrak{E}$ and $\mathfrak{r}$ are system units. Their
definition and values may be found in Appendix \ref{appendix:units}.
$\real{\gamma} = \real{m}_{\parallel}/\real{m}_{\perp}$ is the effective mass
anisotropy.  We may implement $\real{a}$ and $\real{b}$ as unitful or
dimensionless variational parameters. If a dimensionless representation is
desired, we may take $\real{a},\real{b} \rightarrow \mathfrak{r}\real{a},
\mathfrak{r}\real{b}$. When that is done, the $\mathfrak{r}$ drops out of the
equation above and the minimization problem becomes purely numerical.

It can be shown by applying the calculus of variations to quantum mechanics that
an energy eigenstate is an energetic stationary point in state space, and in
particular is a minimum of energy. If selecting from a parameterized subspace,
the problem goes from the calculus of variations to normal real number calculus.
In the \emph{variational method}, we can extend this principle and suggest that
the best approximation to the energy eigenstate consists of selecting a trial
wavefunction from a parameterized set by optimizing the parameters to minimize
the energy. 

Using a value of $\gamma = 0.207901$, we can minimize \eqref{equation:KL-energy}
by varying $\real{a}$ and $\real{b}$. If we do so, we find
\begin{align}
    \real{E}_{\text{min}} &= -1.5653\mathfrak{E},\\
    \real{a}_{\text{min}} &= 0.74832\mathfrak{r},\\
    \real{b}_{\text{min}} &= 0.43001\mathfrak{r}.
\end{align}
This is a ground state energy of $\real{E}_{\text{min}} = \SI{-31.1}{\meV}$.
Since the KL model does not include intervalley coupling, the model predicts a
six-fold degenerate ground state. In Chapter \ref{chapter:results} we will
explore in greater detail how our software package can directly replicate and
freely extend these results. In particular, we wish to examine the implications
of introducing multi-valley interactions.

In their paper, KL suggest a few more trial wavefunctions to approximate various
excited states of the silicon donor. KL found reasonable, but not exact,
agreement with the excited states. We believe that it is worthwhile to consider
the sources of error within the KL model and procedure, especially as one
considers an approach to improve or replace parts of the model.

As an initial examination, we might consider, compare, and contrast the
following two sources of error. To begin with, Let's consider the effective mass
approximation itself. The singular nature of the positive charge center of the
donor is likely to cause some degree of interband coupling, leading to a degree
of divergence between the predictions of the model and the values obtained via
experiment. However, it seems uncertain to this author what degree of violation
of the EMA is induced by the positive charge center. Additionally, the KL model
neglects intervalley contributions to the expectation energy.

Consider the process proceeding from choosing initial assumptions and model
constructions, all the way to error analysis of the final results. We can
imagine that introducing corrections to the EMA model or using a different model
will introduce a change to that process very early on, likely introducing many
downstream effects. If we were wanting to improve our predictions in a timely
manner, we might instead explore a refinement to the prediction process farther
down stream.

As an example of such a downstream error that might be examined is the error
introduced by the use of a variational wavefunction. A variational wavefunction,
in general, is an approximation. The quality of that approximation depends on
many factors, and so cannot be easily refined to produce high-quality
improvements to the wavefunction.

An alternative we can imagine is to express the wavefunction using a functional
basis spanning the Hilbert space of the system in question. Additionally, we can
imagine this basis itself being equipped with variational parameters allowing
for refinement and optimization. Such an approach, combining a basis
representation of the wavefunction with the variational method, is known as the
\emph{Ritz method}.

We can imagine that resolving these two sources of error will each introduce
very different costs. The error that lies farther upstream in the prediction
process will cost more analytical and software development time. In a broader
collaboration, we can also imagine that such an error will require response from
and delegation of work to more parts of the collaboration, introducing
additional time cost in communication and management. An error farther
downstream in the prediction process, on the other hand, will incur far less
time and work costs. We can also imagine that development of such an improvement
could be isolated to smaller working groups of the collaboration, possibly even
a single individual.

Here the authors assert that here we are beginning to see a distinction in
sources of error between what we are calling \emph{technical} and \emph{model}
error. We will explore in more detail what is meant by our usage of these terms,
but as a starting point we can loosely identify that errors lying farther up the
prediction process are \emph{model} errors while those lying father down the
prediction process are \emph{technical} errors.

The exact line between these two categories is contextual and somewhat
arbitrary, and indeed it might be more philosophically accurate to identify the
two as poles in an interpretational spectrum. Regardless of philosophical or
formal ambiguity, we still find the above a useful distinction. As we identified
earlier, this usefulness is particularly pointed when we consider how error
management will be allocated amongst researchers and technicians in a larger
academic or commercial research collaboration. Even the individual researcher
might find themselves in a collaboration of sorts, even if the only members of
such a collaboration are future and past versions of themself. For now, we will
table the discussion of error analysis or improvement of the KL model and
revisit it as part of the historical narrative of this document.



While we can continue examining the other variational wavefunctions employed by
KL in their exploration of the excited donor states of silicon, we find it
convenient at this juncture to move forward. In the next section, we will
examine a significant improvement upon the bespoke single-state variational
method and discuss how R.A. Faulker used the Ritz method.